%% kapa: exam=1 round=35 year=2024 period=1 form=A subject=civil_law subject_ko=민법 seq=17 number=17
%% source: Q-Net 감정평가사 기출문제 (cst00309) artlSeq=5213229 file=raw/1008/5213229/1교시 문제지 원본.pdf
%% conversion: kordoc 4.24.51 -> tools/gen1.py (auto); review=proofread-sonnet; flags=none
\begin{question}{17}
조건과 기한에 관한 설명으로 옳지 않은 것은?
\begin{choices}
  \item 기성조건이 정지조건이면 조건 없는 법률행위가 된다.
  \item 불능조건이 해제조건이면 조건 없는 법률행위가 된다.
  \item 불법조건은 그 조건만이 무효가 되고 그 법률행위는 조건 없는 법률행위로 된다.
  \item 기한은 당사자의 특약에 의해서도 소급효를 인정할 수 없다.
  \item 기한은 원칙적으로 채무자의 이익을 위한 것으로 추정한다.
\end{choices}
\end{question}
\answer{3}
